\lecture{1}

\section{Введение}
Ну короче тыры-пыры все дела там секретно все и тд. 

Сначала поймем, какие модели угроз мы будем рассматривать (какие алгоритмы взломщик может запускать). В основной теории разделяют 2 основных вида:

\begin{enumerate}
    \item Равномерный противник --- полиномиальный вероятностный алгоритм. Причем полином берется не по длине отправленного сообщения, а по параметру \(n\), который называется параметром безопасности (например, для \(512\)-битного шифрования \(n = 512\)).
    \item Неравномерный противник --- набор схем (например, злоумышленник может их физически спаять).
\end{enumerate}

\begin{reminder}[Теорема Эйдлмана]
    \(\mathbf{BPP} \subset \mathbf{P/poly}\)
\end{reminder}

Таким образом, неравномерный противник сильнее равномерного. Заметим, что если \(\mathbf{P} = \mathbf{NP}\), то тогда криптография с открытым ключом не имеет смысла: по открытому ключу можно быстро (за полином) получить закрытый. Однако существуют такие протоколы, которые не требуют неравенства данных двух классов.

В задачах по криптографии обычно Алиса хочет передать сообщение Бобу, а Ева (злоумышленник) пытается взломать шифрование

\begin{example}[Гаммирование (одноразовый блокнот)]
    Пусть Алиса хочет передать \(m\) Бобу. В таком случае они (заранее) договариваются о некотором секрете \(k\), таком, что \(|k| = |m|\). При передаче сообщения, Алиса отправляет \(n = m \oplus k\) (побитово), Боб делает то же самое и получает \(m = n \oplus k\).
\end{example}

Далее мы будем жить в предположении, что \(\mathbf{P} \ne \mathbf{NP}\).

\subsection{Квантовая и постквантовая криптографии}

Относительно недавно выяснилось, что квантовые компьютеры умеют раскладывать числа на множители. Важно понимать, что настоящие \(\mathbf{NPC}\)-задачи (по-видимому, \(\mathsf{FACTORING}\) такой не является) пока что решить при помощи него не получилось. Однако на данный момент вычислительных мощностей даже квантового компьютера не хватает для разложения чисел, используемых в современной криптографии.

Соответственно, постквантовая криптография --- это классические протоколы, устойчивые относительно квантового противника. На данный момент главный претендент на такой протокол --- теория решеток.

Квантовая криптография --- это протокол, использующий квантовый канал связи. Таким образом, если кто-то вдруг будет смотреть за тем, что мы там передаем, то будут помехи и получатель сможет понять, что за каналом следят (далее начинается очень сложная квантовая механика).

Оба данных подхода могут спасти от вмешательства квантового компьютера. 

\section{Односторонние функции}

Интуитивно, односторонняя функиция --- такая, что \(x \mapsto f(x)\) легко вычислить, а \(f(x) \mapsto x\) уже сложно. В силу того, что можно положить \(f(x) = 1\) для всех \(x\), задача обращения односторонней функции звучит так: ''Для данного \(y\) найти хотя бы один \(x\), такой, что \(f(x) = y\)''.

Пусть \(f = \{f_n\}_{n = 1}^\infty\), где \(f_n: \{0, 1\}^{k(n)} \ra \{0, 1\}^{l(n)}\) (регулярная по длине, т.е. при равной длине входа получается равная длина выхода). Здесь \(n\) --- параметр безопасности, \(k, l\) --- полиномы.

\begin{definition}
    \(f\) называется слабо односторонней, если \(\exists q(\cdot)\) --- полином \(\forall R(\cdot)\) --- обратитель (либо равномерный, либо неравномерный противник) \(\exists N \forall n > N: \Pr_{x \in \{0, 1\}^{k(n)}}\{R(f(x)) \in f^{-1}(f(x))\} < 1 - \frac{1}{q(n)}\)
\end{definition}

\begin{definition}
    \(f\) называется сильно односторонней, если \(\forall p(\cdot)\) --- полином \(\forall R(\cdot)\) --- обратитель (либо равномерный, либо неравномерный противник) \(\exists N \forall n > N: \Pr_{x \in \{0, 1\}^{k(n)}}\{R(f(x)) \in f^{-1}(f(x))\} < \frac{1}{q(n)}\)
\end{definition}

\begin{note}
    На самом деле, выше написаны 4 определения --- по два для равномерного и неравномерного противника
\end{note}

\begin{note}
    Предположительно, функция \(f: \{0, 1\}^n \times \{0, 1\}^n \ra \{0, 1\}^{2n}, f(x, y) = x \cdot y\) является односторонней.
\end{note}

\begin{theorem}[Об усилении]
    Существует слабо односторонняя функция \(\Ra\) существует сильно односторонняя функция.
\end{theorem}

Для доказательства теоремы потребуются несколько утверждений:

\begin{proposition}
    Пусть \(f\) --- сильно односторонняя функция. Пусть \(g(x, y) = (f(x), y)\). Тогда \(g\) --- сильно односторонняя.
\end{proposition}
\begin{proof}
    Очевидно, что \(g\) --- регулярная по длине полиномиальная функция. Пусть теперь у \(g\) есть обратитель \(R_g\). Покажем, что тогда и у \(f\) есть обратитель \(R_f\). Положим \(R_f(z) = (R_g(z, y))_x\). Далее нетрудно можно посчитать распределение полученных ответов и доказать утверждение.
\end{proof}

\begin{proof}[Конструкция теоремы об усилении]
    Рассмотрим функцию \(F(x_1, \dots x_M) = (f(x_1), \dots f(x_M))\) (здесь \(M = M(n)\)). Если обратитель будет наивным и будет обращать вектор покомпонентно, то он угадает с вероятностью \(\left( 1 - \frac{1}{q(n)} \right)^M\), что можно сделать сколь угодно малым.
\end{proof}

\lecture{2}

Однако, обратитель вовсе не обязан быть наивным --- он может действовать как-то по-другому, необязательно обращая функцию покомпонентно. Для доказательства, мы не будем из обратителя функции \(f\) строить обратитель \(F\). Наоборот, из обратителя \(F\) мы построим очень хороший обратитель \(f\).

\begin{proof}[Доказательство теоремы об усилении]
    Предположим, что \(F\) не сильно односторонняя. Тогда 
    \[\exists q(\cdot) \exists R_f \forall N \exists n > N \Pr\{R(F(x)) \in F^{-1}(F(x))\} > \frac{1}{q(n)}\]
    Построим по \(R_F\) обратитель \(R_f\) функции \(f\). Пусть мы хотим обратить \(f^{-1}(y) = \dots\). Дополним \(y\) до \((y, f(x_2), \dots f(x_M))\) (заметим, что после такого дополнения полученный вектор имеет такое же распределение, как и \(x\) из определения \(R_F\)) и рассмотрим \((\vec{z_1}, \dots \vec{z_M}) = R_F(y, f(x_2), \dots f(x_M))\), после чего проверим \(f(z_1) = y\). Повторим данную операцию \(Q\) раз, причем сделаем эти \(Q\) итераций для каждого места, где может стоять \(y\) (проще говоря, постваим \(y\) на всевозможные места и для каждого из них проведем данную операцию независимо \(Q\) раз). Если мы увидим хотя бы один успех, то возвращаем его. Так строится \(R_f\). Докажем теперь, что он действительно быстро обращает функцию \(f\).

    Рассмотрим функцию \(S_i(x) = \Pr\bigg\{\big(R_F(f(x_1), \dots f(x_{i - 1}), f(x), f(x_{i + 1}), \dots f(x_M))\big)_i \in f^{-1}(f(x))\bigg\}\) и положим \(S(x) = \max \{S_1(x), \dots S_M(x)\}\). Выберем некоторый параметр \(\epsilon \sim \frac{1}{poly(n)}\) (определим его позже) и разобьем слова на простые: для которых \(S(x) \ge \epsilon\) и сложные: для которых \(S(x) < \epsilon\). Также положим \(\delta\) --- доля сложных слов.

    \begin{lemma}
        \(\Pr \{R_f\text{ ошибся}\} \le \delta + (1 - \delta)(1 - \epsilon)^Q\)
    \end{lemma}
    \begin{proof}
        \[\Pr \{R_f\text{ ошибся}\} \le \underbrace{\Pr\{x\text{ сложное}\}}_\delta + \underbrace{\Pr\{x\text{ простое}\}}_{1 - \delta} \cdot \underbrace{\Pr\{R_f\text{ ошибся} | x\text{ простое}\}}_{(1 - \epsilon)^Q}\]
    \end{proof}

    \begin{lemma}
        \(\Pr\{R_F\text{ успешен}\} \le M\delta\epsilon + (1 - \delta)^M\)
    \end{lemma}
    \begin{proof}
        Покроем все вероятностное пространство событиями:
        \[\begin{array}{c|c}
            \text{Событие} & \Pr \\
            \hline
            \text{Все \(x_i\) --- простые} & (1 - \delta)^M \\
            \hline
            \text{\(x_1\) --- простые} & \delta \\
            \hline
            \vdots & \vdots \\
            \hline
            \text{\(x_M\) --- простые} & \delta \\
        \end{array}\]
        Соответственно, по формуле полной вероятности:
        \[\Pr\{R_F\text{ успешен}\} \le \Pr\{R_F\text{ успешен} | \text{все \(x_i\) простые}\}\cdot (1 - \delta)^M +\]
        \[+ \sum_{i = 1}^M \Pr\{R_F\text{ успешен} | \text{\(x_i\) сложное}\}\cdot \delta \le M\delta\epsilon + (1 - \delta)^M\]
    \end{proof}

    Поскольку мы предположили противное, то \(\Pr\{R_f\text{ ошибается}\} \ge \frac{1}{p(n)}\) и \(\Pr\{R_F\text{ успешен}\} \ge \frac{1}{q(n)}\). Подберем параметры \(M, Q, \epsilon\) так, чтобы получилось противоречие с доказанными леммами:

    \begin{enumerate}
        \item \(Q = \frac{n}{\epsilon} \Ra (1 - \epsilon)^Q \approx e^{-n} \Ra \delta \ge \frac{1}{2p(n)}\)
        \item \(M = np(n) \Ra (1 - \delta)^M \sim e^{-\frac{n}{2}}\Ra M\delta\epsilon \ge \frac{1}{2q(n)} \Ra \epsilon \ge \frac{1}{2q(n)M\delta} \ge \frac{1}{2q(n)M}\)
        \item Наконец, положим \(\epsilon = \frac{1}{4q(n)np(n)}\) --- получим противоречие с неравенством выше.
    \end{enumerate}
\end{proof}

\subsection{Односторонние перестановки}

Рассмотрим одностороннюю функцию, являющуюся биекцией, т.е. \(f = \{f_n\}_{n = 1}^\infty\), где \(f_n: \{0, 1\}^{k(n)} \ra \{0, 1\}^{k(n)}\).

К сожалению, на данный момент нет даже кандидатов в примеры таких функций (на всём \(\{0, 1\}^{k(n)}\)). В связи с этим рассматриваются более слабые свойства: биекция \(f_n: D_n \ra D_n\), где \(D_n \subset \{0, 1\}^n\). То есть мы разрешаем функции быть перестановкой не на всех строках длины \(n\), а только на некотором подмножестве.

\begin{example}[Кандидат в примеры]
    Рассмотрим функцию Рабина: \((m, x) \mapsto (m, x^2 \mod m)\), где \(x \underset{m}{\equiv} y^2\). Для \(m = 4k + 3\) --- простого данная функция будет легко обратимой: достаточно взять \(R_f(z) = z^{\frac{m + 1}{4}}\). Тогда: \(R_f(x^2) = R_f(y^4) = y^2 = x\).
    
    Теперь рассмотрим \(m = pq\), где \(p, q \underset{4}{\equiv} 3\). Если \(p, q\) известны, то обратить функцию легко алгоритмом выше, после чего использовать КТО. В случае же, когда \(p, q\) неизвестны, обратить её уже становится сложнее.
\end{example}

Заметим, что в примере выше область определения --- это не все строки длины \(n\), а именно некоторое подмножество \(D_n\). Но тут возникает естественный вопрос: а как нам вообще работать с элементами из \(D_n\)? Если мы хотим, скажем, выбрать случайный элемент из \(D_n\) (чтобы подать его на вход функции), то нам нужно уметь его эффективно сгенерировать. Причём сгенерировать так, чтобы распределение было ``почти равномерным'' на \(D_n\) --- иначе противник может использовать неравномерность в свою пользу.

\begin{definition}
    В общем определении требуется, чтобы за полиномиальное время можно было сгенерировать распределение, пренебрежимо мало отличимое от равномерного на \(D_n\). Более формально, пусть \(\alpha\) --- сгенерированное распределение, тогда
    \[\forall p(\cdot) \exists N: \forall n > N \forall S \subset D_n \left| \alpha(S) - \frac{|S|}{|D_n|}\right| < \frac{1}{p(n)}\]
\end{definition}

\lecture{3}

\subsection{Односторонние перестановки с секретом}

\textit{Примечание от лектора:} по-другому: ''one-way trapdoor permutations''

Обычные односторонние перестановки функция --- те, для которых \(f_\alpha(x)\) легко вычислима, а вот \(f_\alpha^{-1}(y)\) --- трудно (если знать только \(\alpha\)). Перестановка с секретом --- та, для которой \(f_\alpha^{-1}(y)\) легко вычисляется, если знать \(\alpha\) и некоторый секрет \(\tau\). Дадим этому формальное определение:

\begin{definition}
    Односторонняя перестановка с секретом --- функция \(f: \{0, 1\}^{k(n)} \times \{0, 1\}^{l(n)} \ra \{0, 1\}^{l(n)}\) (мы будем обозначать \(f_\alpha(x) := f(\alpha, x)\)), для которой существует множества \(D_\alpha \subset \{0, 1\}^{l(n)}\), причем \(f: D_\alpha \ra D_\alpha\) --- биекция и 4 полиномиальных вероятностных алгоритма:
    \begin{enumerate}
        \item Генератор (генерирует ключи) \(G: \text{случаные биты} \mapsto \text{пара}(\alpha, \tau)\) 
        \item Сэмплер (сэмплирует \(D_\alpha\)) \(S: \alpha \mapsto \text{почти случайный элемент } D_\alpha\)
        \item Вычислитель функции \(F: (\alpha, x) \mapsto f_\alpha(x)\)
        \item Обратитель \(B: (\alpha, \tau, y) \mapsto f_\alpha^{-1}(y)\)
    \end{enumerate}
    Кроме этих алгоритмов, еще требуется односторонность перестановки:
    \[\forall p \forall R \exists N \forall n > N: \Pr_{\alpha, y}\{R(\alpha, y) = f^{-1}_\alpha(y)\} < \frac{1}{p(n)}\]
\end{definition}

\begin{example}[Предположительный пример]
    Рассмотрим функцию Рабина, где \(\alpha = p \cdot q\) (\(p, q\) --- простые вида \(4k + 3\)) --- модуль, по которому производим арифметику, \(D_\alpha\) --- квадратичные вычеты \(\mod \alpha\), секрет --- \(\tau = (p, q)\), \(f_\alpha(x) = x^2 \mod \alpha\).
\end{example}

\begin{note}
    Про данную функцию неизвестно, действительно ли она является необратимой, но точно можно сказать, что квантовым компьютером данный алгоритм можно взломать (а именно, разложить на множители число \(pq\)). Таким образом, данная функция является примером \textbf{НЕ} постквантовой криптографии.
\end{note}

\section{Генераторы случайных чисел}
Можно генерировать числа по-разному. Можно рассматривать ГИСЧ --- генераторы истинно случайных чисел и ГПСЧ --- генераторы псевдо-случайных чисел. Последние почти всегда работают по следующему правилу: дается некоторый \(x\) малой длины, по которому строится \(G(x)\) большой длины. Нетрудно понять, что таким образом получить равномерное распределение на \(\{0, 1\}^{|G(x)|}\) нельзя (т.к. \(|G(x)| > |x|\)), но можно попытаться получить нечто ''похожее'' на случаное распределение.

\begin{definition}
    Статистическое расстояние между распределениями случайных величин \(\alpha_n, \beta_n\) определяется как
    \[\operatorname{dist}(\alpha_n, \beta_n) = \max_{S \subset \{0, 1\}^n} |\Pr\{\alpha_n \in S\} - \Pr\{\beta_n \in S\}|\]
\end{definition}

\begin{proposition}
    \[\operatorname{dist}(\alpha_n, \beta_n) = \frac{1}{2}\sum_{x \in \{0, 1\}^n} |\Pr\{\alpha_n = x\} - \Pr\{\beta_n = x\}|\]
\end{proposition}
\begin{proof}
    Пусть максимум в определении достигается на множестве \(S\). Тогда он достигается на множестве \(\overline{S} = \{0, 1\}^n \setminus S\). Б.О.О, \(\Pr\{\alpha_n \in S\} \ge \Pr\{\beta_n \in S\}\). Тогда \(\forall x \in S: \Pr\{\alpha_n = x\} \ge \Pr\{\beta_n = x\}\) (иначе можно ''убрать'' невыгодный \(x\) и значение только увеличится). Аналогично \(\forall x \in \overline{S}: \Pr\{\alpha_n = x\} \le \Pr\{\beta_n = x\}\). Заметим тогда, что
    \[\Pr\{\alpha_n \in S\} - \Pr\{\beta_n \in S\} = \sum_{x \in S} \Pr\{\alpha_n = x\} - \Pr\{\beta_n = x\} = \sum_{x \in S} |\Pr\{\alpha_n = x\} - \Pr\{\beta_n = x\}|\]
    Аналогично:
    \[\Pr\{\alpha_n \in \overline{S}\} - \Pr\{\beta_n \in \overline{S}\} = \sum_{x \in S} \Pr\{\beta_n = x\} - \Pr\{\alpha_n = x\} = \sum_{x \in \overline{S}} |\Pr\{\alpha_n = x\} - \Pr\{\beta_n = x\}|\]
    Но тогда:
    \[\operatorname{dist}(\alpha_n, \beta_n) = \frac{\left( \Pr\{\alpha_n \in S\} - \Pr\{\beta_n \in S\} \right) + \left( \Pr\{\beta_n \in \overline{S}\} - \Pr\{\alpha_n \in \overline{S}\} \right)}{2} = \]
    \[= \frac{1}{2}\sum_{x \in \{0, 1\}^n} |\Pr\{\alpha_n = x\} - \Pr\{\beta_n = x\}|\]
\end{proof}

\begin{note}
    Таким образом, \(\operatorname{dist}\) является метрикой (\(L_1\)-метрикой) на множестве всех случайных величин.
\end{note}

\begin{definition}
    Будем говорить, что \(\alpha_n, \beta_n\) статистически близки (\(\alpha_n \sim \beta_n\)), если 
    \[\forall p(\cdot) \exists N \forall n > N: \operatorname{dist}(\alpha_n, \beta_n) < \frac{1}{p(n)}\]
\end{definition}

\begin{note}
    Из утверждения выше получается эквивалентное определение:
    \[\alpha_n \sim \beta_n \Lra \forall p(\cdot) \forall T \exists N \forall n > N: |\Pr\{T(\alpha_n) = 1\} - \Pr\{T(\beta_n) = 1\}| < \frac{1}{p(n)}\]
    Здесь \(T\) --- любая булева функция
\end{note}

\begin{definition}
    Будем говорить, что \(\alpha_n, \beta_n\) вычислительно неотличимы (\(\alpha_n \sim \beta_n\), будет понятно из контекста, что имеется в виду), если 
    \[\forall p(\cdot) \forall T \exists N \forall n > N: |\Pr\{T(\alpha_n) = 1\} - \Pr\{T(\beta_n) = 1\}| < \frac{1}{p(n)}\]
    Где \(T\) --- противник (полиномиально вычислимый вероятностный алгоритм или семейство схем полиномиальной длины).
\end{definition}

\begin{note}
    Статистическая близость и вычислительная неотличимость --- отношения эквивалентности
\end{note}
\begin{proof}
    Рефлексивность и симметричность очевидны. Докажем транзитивность: рассмотрим моменты \(N_1, N_2\), для которых \(\forall n > N_1, N_2: \operatorname{dist}(\alpha_n, \beta_n) < \frac{1}{2p(n)}, \operatorname{dist}(\beta_n, \gamma_n) < \frac{1}{2p(n)}\). Но тогда \(\forall n > \max\{N_1, N_2\}: \operatorname{dist}(\alpha_n, \gamma_n) < \frac{1}{p(n)}\)
\end{proof}

\begin{proposition}
    Пусть \(\alpha_n \sim \beta_n\) (вычислительно неотличимы) и \(f\) --- полиномиально вычислимая функция. Тогда \(f(\alpha_n) \sim f(\beta_n)\) --- тоже вычислимо неотличимы.
\end{proposition}
\begin{proof}
    Пусть \(\exists p \exists T \forall N \exists n > N: |\Pr\{T(f(\alpha_n)) = 1\} - \Pr\{T(f(\beta_n)) = 1\}| \ge \frac{1}{p(n)}\).
    Рассмотрим \(T' = T \circ f\) --- полиномиально вычислимый. Но тогда  \(\forall N \exists n > N: |\Pr\{T'(\alpha_n) = 1\} - \Pr\{T'(\beta_n) = 1\}| \ge \frac{1}{p(n)}\), получили, что \(\alpha_n \not\sim \beta_n\).
\end{proof}

\begin{definition}
    Генератор псевдо-случайных чисел --- полиномиально вычислимая функция (детерминированным алгоритмом) \(G_n: \{0, 1\}^n \ra \{0, 1\}^{p(n)}, p(n) > n\), такая, что \(G(x) \sim y\), где \(y\) --- случайные биты длины \(p(n)\) (при условии, что \(|x| = n\)). Противник для ГПСЧ обычно называют отличителем.
\end{definition}

\begin{note}
    Обычно конструкция ГПСЧ выглядит следующим образом: имеется \(g: \{0, 1\}^n \ra \{0, 1\}^{n + 1}\). Соответствено \(g(x) = (y, b)\) (\(b\) --- последний бит). Соответственно, вывод определяется следующим образом: выдаем последний бит и передаем \(y\) обратно на вход функции \(g\), так делаем сколько угодно раз.
\end{note}

\begin{theorem}
    \(\exists\) ГПСЧ \(\Ra \exists\) односторонняя функция
\end{theorem}
\begin{proof}
    Докажем, что сам генератор будет являться односторонней функцией. Пусть \(R\) --- обратитель для \(G\). Тогда для бесконечно многих \(n\): \(\Pr_{x \in \{0, 1\}^n}\{R(G(x)) \in G^{-1}(G(x))\} = \epsilon(n) \ge \frac{1}{q(n)}\). Рассмотрим функцию
    \[T(y) = \left\{\begin{array}{l}
        1, G(R(y)) = y \\
        0, \text{иначе}
    \end{array}\right.\]
    Заметим, что \(|G(\{0, 1\}^n)| \le 2^{n - 1}\). В таком случае:
    \[\Pr\{T(y) = 1\} = \left\{\begin{array}{l}
        \epsilon, y\text{ из распределения }G(x) \\
        \frac{|G(\{0, 1\}^n)|}{2^n} \le \frac{\epsilon}{2}, \text{\(y\) --- случайный}
    \end{array}\right.\]
    Но тогда \(\big|\Pr_{y \in \{0, 1\}^{p(n)}}\{T(y) = 1\} - \Pr_{x \in \{0, 1\}^{n}}\{T(G(x)) = 1\}\big| \ge \frac{\epsilon(n)}{2} \ge \frac{1}{2q(n)}\), поэтому \(G(x) \not\sim y\), получили противоречие.
\end{proof}

\lecture{4}

В книжке Верещагина и Шеня можно найти доказательство следующей сложной теореме:

\begin{theorem}[б/д]
    \(\exists\) односторонняя функция \(\Ra\) существует ГПСЧ
\end{theorem}

Мы докажем следующее более слабое утверждение:

\subsection{Теорема о существовании ГПСЧ из односторонней функции}

В данном разделе будем под односторонней перестановкой понимать перестановку с областью определения \(\{0, 1\}^n\)

\begin{definition}
    Пусть \(h: \{0, 1\}^n \ra \{0, 1\}, g: \{0, 1\}^n \ra \{0, 1\}^n\) --- полиномиально вычислимые функции. Функция \(h\) называется трудным битом для функции \(g\), если
    \[\forall q \forall P \exists N \forall n > N \left|\Pr\{P(g(x)) = h(x)\} - \frac{1}{2}\right| \le \frac{1}{q(n)}\]
    \(P\) обычно называют предсказателем
\end{definition}

\begin{note}
    В частности, функция \(g\) должна быть труднообратимой.
\end{note}

\begin{theorem}
    \(\exists\) односторонняя перестановка \(\Ra \forall p(n) \exists\) ГПСЧ \(G: \{0, 1\}^n \ra \{0, 1\}^{p(n)}\).
\end{theorem}
\begin{proof}[Схема доказательства]
    Пусть \(f: \{0, 1\}^n \ra \{0, 1\}^n\) --- односторонняя перестановка. По ней мы получим односторонюю перестановку с трудным битом, а именно: \(g(x, y) = (f(x), y)\), а \(h(x, y) = x \odot y = \sum_{i = 1}^n x_iy_i(\operatorname{mod} 2)\). После этого получим \(G(x) = g(x)h(x)\), и общий случай \(G'(x) = h(x)h(g(x))\dots h(g^{p(n) - 1}(x))\) --- что и будет являться ГПСЧ для данного полинома.
\end{proof}

\begin{note}
    По схеме доказательства можно заметить следующее: если мы прочитаем у первые \(n\) бит выхода построенного ГПСЧ \(G\), то не сможем (это будет трудно) вычислить оставшийся бит.
\end{note}

Докажем каждый переход в теореме:

\begin{lemma}[XOR-Лемма]
    Пусть \(g\) --- односторонняя перестановка с трудным битом \(h\). Тогда \(G(x) = g(x)h(x)\) является ГПСЧ
\end{lemma}
\begin{proof}
    Пусть \(D\) --- отличитель для \(G\). Покажем, как по нему построить \(P\) --- предсказатель для трудного бита \(h\). По определению отличителя (запишем явно отрицание к вычислительной неотличимости):
    \[\exists q \forall N \exists n > N |\Pr\{D(g(x)h(x)) = 1\} - \Pr\{D(yr) = 1\}| > \frac{1}{q(n)}\]
    Заметим, что \(g\) --- односторонняя перестановка с областью определения \(\{0, 1\}^n\), поэтому \(x \sim \mathcal{U}(\{0, 1\}^n) \Ra g(x) \sim \mathcal{U}(\{0, 1\}^n)\). Поэтому данное условие можно переписать так:
    \[\exists q \forall N \exists n > N |\Pr\{D(g(x)h(x)) = 1\} - \Pr\{D(g(x)r) = 1\}| > \frac{1}{q(n)}\]
    Положим \(P(y) = D(yr) \oplus r \oplus 1\), где \(r\) --- случайный бит. Тогда значения \(P\) определяются по следующей таблице:
    \[\begin{array}{c|c|c}
        \substack{D(g(x)1) \ra\\D(g(x)0) \downarrow} & 0 & 1 \\
        \hline
        0 & \text{случайно} & 1 \\
        \hline
        1 & 0 & \text{случайно}\\
    \end{array}\]
    Также положим \(\alpha, \beta, \gamma, \delta\) --- вероятности попадения в клетки, т.е.:
    \[\begin{array}{c|c|c}
        \substack{D(g(x)1) \ra\\D(g(x)0) \downarrow} & 0 & 1 \\
        \hline
        0 & \alpha & \beta \\
        \hline
        1 & \gamma & \delta\\
    \end{array}\]
    Обозначим \(\alpha = \alpha_0 + \alpha_1\) --- вероятности событий \(h(x) = 0, h(x) = 1\) внутри соответствющей клетки (например, \(\beta_0 = P(h(x) = 0, D(g(x)0) = 0, D(g(x)1) = 1)\)).
    Заметим, что:
    \[\Pr\{D(g(x)h(x)) = 1\} = \delta_0 + \delta_1 + \beta_1 + \gamma_0\]
    \[\Pr\{D(yr) = 1\} = \delta_0 + \delta_1 + \frac{\beta_0 + \beta_1}{2} + \frac{\gamma_0 + \gamma_1}{2}\]
    Соответственно:
    \[\Pr\{D(g(x)h(x)) = 1\} - \Pr\{D(yr) = 1\} = \delta_0 + \delta_1 + \beta_1 + \gamma_0 - \delta_0 - \delta_1 - \frac{\beta_0 + \beta_1}{2} - \frac{\gamma_0 + \gamma_1}{2} =\]
    \[= \frac{\beta_1 - \beta_0}{2} + \frac{\gamma_1 - \gamma_0}{2} > \epsilon\]
    Тогда:
    \[\Pr\{P(g(x)) = h(x)\} = \frac{\alpha_0 + \alpha_1}{2} + \frac{\delta_0 + \delta_1}{2} + \beta_1 + \gamma_0 =\]
    \[\frac{\alpha_0 + \alpha_1}{2} + \frac{\beta_0 + \beta_1}{2} + \frac{\gamma_0 + \gamma_1}{2} + \frac{\delta_0 + \delta_1}{2} + \frac{\beta_1 - \beta_0}{2} + \frac{\gamma_1 - \gamma_0}{2} > \frac{1}{2} + \epsilon\]
    Получили, что \(h\) --- не случайный бит для \(g\).
\end{proof}

\begin{lemma}
    Пусть есть ГПСЧ: \(G: \{0, 1\}^{n} \ra \{0, 1\}^{n + 1}\). Тогда существует ГПСЧ: \(G': \{0, 1\}^{n} \ra \{0, 1\}^{p(n)}\).
\end{lemma}
\begin{proof}
    Положим \(G(x) = g(x)h(x)\), где \(g: \{0, 1\}^{n} \ra \{0, 1\}^{n}, h: \{0, 1\}^{n} \ra \{0, 1\}\) --- полиномиально вычислимые функции (проще говоря, \(h\) --- последний бит \(G\), а \(g\) --- его префикс). Определим \(G_2(x) = G(G(x))h(x) = g(g(x))h(g(x))h(x)\). Положим \(G_0(zr_1r_2) = zr_1r_2, G_1(yr) = g(y)h(y)r\). Заметим, что \(G_0 \sim G_1\) (в смысле вычислительной неотличимости), так как \(G\) --- генератор, т.е. \(g(y)h(y) \sim zr_1\). Рассмотрим преобразование \(F(zr_1r_2) = g(z)h(z)r_1\). Заметим, что \(G_1(yr) = F(G_0(yrr_2))\), а \(G_2(x) = F(G_1(xr))\). Так как \(F\) полиномиально вычислима, то \(G_1 \sim G_2\), получаем \(G_0 \sim G_2\), т.е. \(G_2\) --- ГПСЧ.

    % Осталось показать, что построенный \(G_k\) действительно является ГПСЧ. Положим \(k = p(n) - n\). Для \(i = 0, \dots, k\) рассмотрим
    % \[
    % G_i(y, r_{i+1}, \dots, r_k) = g^i(y)h(g^{i-1}(y))\dots h(y)r_{i+1}\dots r_k,
    % \]
    % где \(|y| = n\), а \(r_j\) --- отдельные случайные биты. Тогда \(G_0\) --- тождественная функция на \(n+k\) битах, а
    % \[
    % G_k(x) = g^k(x)h(g^{k-1}(x))\dots h(x).
    % \]
    % Пусть \(H_i = G_i(\text{случайные биты длины } n+k-i)\). Тогда \(H_0\) --- равномерное распределение на \(\{0,1\}^{n+k}\), а \(H_k = G_k(x)\) при случайном \(x\).

    % Покажем, что \(H_i \sim H_{i+1}\). Действительно, если \(z\) --- случайные \(n\) бит, то
    % \[
    % H_{i+1} = G_i(g(z), h(z), r_{i+2}, \dots, r_k).
    % \]
    % Но \(g(z)h(z) = G(z) \sim \text{случайные } n+1 \text{ бит}\), поэтому по утверждению о сохранении вычислительной неотличимости при полиномиально вычислимых преобразованиях
    % \[
    % H_{i+1} \sim G_i(\text{случайные } n+1 \text{ бит}, r_{i+2}, \dots, r_k) = H_i.
    % \]
    % Значит, все соседние \(H_i\) и \(H_{i+1}\) вычислительно неотличимы. По транзитивности \(H_0 \sim H_k\), то есть \(G_k(x)\) вычислительно неотличимо от равномерного распределения на \(\{0,1\}^{n+k}\). Поэтому можно положить \(G'(x) = G_k(x)\).

    % Если противник неравномерный и различает \(H_0\) и \(H_k\) с преимуществом \(\epsilon\), то по неравенству треугольника найдётся \(i\), для которого он различает \(H_i\) и \(H_{i+1}\) с преимуществом хотя бы \(\epsilon/k\). Это \(i\) зашиваем в схему-отличитель для исходного \(G\); для равномерного противника \(i\) можно выбрать случайно.

    % Таким образом можно построить \(G_k(x): \{0, 1\}^n \ra \{0, 1\}^{n + k}\), где \(k \in \N\), который будет определяться как \(G_k(x) = g^{k}(x)h(g^{k - 1}(x))\dots h(g(x))h(x)\) и будет являться ГПСЧ.

    % Теперь положим \(m(n) = p(n) - n\) и введем \(G_{m(n)} = g^{m(n)}(x)h(g^{m(n) - 1}(x))\dots h(g(x))h(x)\) --- полиномиально вычислимая функция. Докажем, что она будет ГПСЧ.

    % Предположим противное, тогда \(G_0 \not\sim G_{m(n)}\), где \(G_0\) --- равномерное распределение на \(\{0, 1\}^{m(n)}\). Тогда:
    % \[\exists q \exists D \forall N \exists n > N |\Pr\{D(G_{m(n)}(x)) = 1\} - \Pr\{D(G_0(y)) = 1\}| > \frac{1}{p(n)}\]
    % По неравенству треугольника:
    % \[\sum_{i = 1}^{m(n)}|\Pr\{D(H_{i + 1}(x)) = 1\} - \Pr\{D(H_{i}(y)) = 1\}| \ge |\Pr\{D(G_{m(n)}(x)) = 1\} - \Pr\{D(G_0(y)) = 1\}| > \frac{1}{p(n)}\]
    % Поэтому 
    % \[\exists q \exists D \forall N \exists i \exists n > N |\Pr\{D(G_{i + 1}(x)) = 1\} - \Pr\{D(G_i(y)) = 1\}| > \frac{1}{m(n)p(n)}\]

    %     Зафиксируем это \(i\). Так как противник неравномерный, зашьём \(i\) в схему. Рассмотрим
    % \[
    %     D_i = D \circ F^i.
    % \]
    % Напомним, что \(H_i = G_i(x)\|U_{m(n)-i}\), причём \(F(H_i) = H_{i+1}\). Поэтому
    % \[
    %     D_i(H_0) = D(H_i), \qquad D_i(H_1) = D(H_{i+1}).
    % \]
    % Следовательно,
    % \[
    %     \bigl|\Pr\{D_i(H_0) = 1\} - \Pr\{D_i(H_1) = 1\}\bigr|
    %     >
    %     \frac{1}{m(n)q(n)}.
    % \]
    % Но \(H_0 = U_{n+m(n)}\), а \(H_1 = G(x)\|U_{m(n)-1}\). Значит, \(D_i\) различает
    % \(U_{n+m(n)}\) и \(G(x)\|U_{m(n)-1}\) с преимуществом больше \(1/(m(n)q(n))\).

    % Зашьём теперь те случайные биты \(r \in \{0,1\}^{m(n)-1}\), на которых достигается максимум преимущества, и положим
    % \[
    %     D'(w) = D_i(w\|r).
    % \]
    % Тогда
    % \[
    %     \bigl|\Pr\{D'(U_{n+1}) = 1\} - \Pr\{D'(G(x)) = 1\}\bigr|
    %     >
    %     \frac{1}{m(n)q(n)}.
    % \]
    % Но \(m(n) = p(n) - n\) --- полином, поэтому правая часть не является пренебрежимо малой. Получили противоречие с тем, что \(G\) --- ГПСЧ. Значит, \(G_{m(n)}\) действительно является ГПСЧ.

    % В случае равномерного противника вместо зашивания \(i\) можно угадать его с вероятностью \(1/m(n)\). Тогда преимущество уменьшится до \(1/(m(n)^2q(n))\), что по-прежнему не пренебрежимо мало. Поэтому лемма верна и для равномерного противника.

    Для произвольного \(G: \{0, 1\}^n \ra \{0, 1\}^{p(n)}\) пока не написал доказательство.
\end{proof}
